\chapter*{Apresentação}
\addcontentsline{toc}{chapter}{Apresentação}

Este caderno reúne fichas de leitura da literatura sobre o padrão-ouro selecionada para o capítulo 2 da dissertação. A seleção seguiu um critério de compatibilidade com a pergunta do capítulo: situar o debate brasileiro sobre a Caixa de Conversão no interior da discussão internacional sobre o padrão monetário, identificando as vertentes de pensamento, os agentes e o vocabulário que os jornais de 1906 a 1914 herdaram ou rejeitaram.

As fichas seguem um modelo fixo. Cada uma registra a referência completa, o problema e a tese do texto, o argumento em seus passos principais, os conceitos com a definição operacional adotada pelo autor, a evidência e o método, os agentes e vertentes nomeados, as pontes para as fontes primárias, as críticas e limites, e o uso previsto no capítulo. As citações diretas são curtas e sempre acompanhadas da página.

A ordem das partes antecipa a estrutura do capítulo: o debate internacional, o papel dos banqueiros e do crédito soberano, a importação das categorias para o debate brasileiro e a comparação com a Caja de Conversión argentina. A síntese final organiza o que as fichas oferecem a cada parte.
